\label{sec:legal}

\subsection{Neighborhood Stability as a Community of Interest}

The communities of interest doctrine, as codified in state redistricting
criteria from California to Colorado to Arizona, does not enumerate the specific
types of community that qualify~\citep{ncsl2021criteria}.
California's definition---a ``contiguous population which shares common social
and economic interests that would be best served by common representation''
(Cal.\ Gov.\ Code \S 8252(c))---is functional, not categorical.
Housing character satisfies this functional definition on multiple independent
grounds.

\paragraph{Homeowner financial stake.}
A census tract in which 75\% of units are owner-occupied is a population with
a direct and legally enforceable financial stake in federal housing finance
policy, historic tax credits, Community Development Block Grant allocations,
and infrastructure maintenance.
These federal programs affect homeowner equity directly.
Renters in an adjacent apartment corridor have no such equity stake and may
have opposite preferences on, for example, housing voucher policy or
inclusionary zoning mandates.
The owner-occupancy rate is therefore a legally cognizable marker of shared
political interest, distinct from and complementary to the economic character
markers in M.1.

\paragraph{Long-tenured homeowner representation.}
Courts reviewing redistricting plans have accepted testimony about the
representation interests of long-established communities as evidence of
community of interest.
In \textit{Harris v.\ Arizona Independent Redistricting Commission},
136 S.Ct.\ 1301 (2016)~\citep{harris2016}, the Supreme Court upheld the
Commission's plan against a partisan gerrymandering claim in part because
the Commission had documented genuine community bonds among the affected
populations.
The housing vintage score operationalizes the long-tenured homeowner community:
a tract with median year built before 1960 and high owner-occupancy is,
by statistical inference, a community with residents who have owned their
homes for decades.
Their representation interests on matters of historic preservation, estate
taxation, and infrastructure maintenance are shared and legally recognizable.

\subsection{Historic Preservation Districts}

The National Historic Preservation Act of 1966 (NHPA) established the
framework for federal recognition of historic districts, codified at 16 U.S.C.\
\S 470~\citep{nhpa1966}.
State enabling legislation in all 50 states implements parallel historic
preservation programs.
North Carolina's State Historic Preservation Office maintains the National
Register of Historic Places nominations for neighborhoods including Cameron Park
in Raleigh, Trinity Park in Durham, and Dilworth in Charlotte.

The relevance for redistricting is twofold.
First, historic district designation is an administrative recognition that
a neighborhood's built environment constitutes a community bond---precisely
the kind of administrative zone co-membership that M.6 operationalizes.
Second, and more directly relevant to M.3, the housing vintage score captures
the underlying physical characteristic (age of construction) that qualifies
a neighborhood for historic district designation.
A redistricting plan that splits a pre-1940 neighborhood across district
boundaries may disrupt the coordinated advocacy of historic preservation
advocates for federal programs that require unified district representation.

\paragraph{Oregon and Colorado precedents.}
Oregon's 2011 redistricting criteria, as applied by the Legislative Assembly,
accepted expert testimony citing historic neighborhood cohesion as a basis
for keeping Portland's Northeast neighborhood together in a single district.
Colorado's Independent Legislative Redistricting Commission in 2021 explicitly
identified ``historic communities'' as a recognized sub-type of community of
interest under Article V, Section 44 of the state constitution~\citep{coart5s44}.
The housing vintage score provides a quantitative operationalization of the
``historic community'' designation that is legally grounded and publicly verifiable.

\subsection{Property Tax Fiscal Bond in the Housing Context}

The property tax fiscal bond (M.0 and M.6) is most directly relevant to
housing character in the school district context.
Single-family homeowner neighborhoods are the primary funding base for
local public schools: property tax revenue from owner-occupied single-family
homes constitutes the majority of school district revenue in most states.
A tract with high SF and high OO is a tract that is invested in local school
quality in a way that a high-rental, high-MF tract is not.

Colorado's redistricting criteria recognize this bond by listing ``economic
interests'' and ``trade area'' community membership alongside traditional
communities of interest~\citep{coart5s44}.
The homevoter hypothesis~\citep{fischel2001} formalizes the connection:
homeowners in a school district are co-investors in a fiscal commons that
gives them shared political interests distinct from those of renters in
adjacent rental-heavy tracts.

\subsection{Partisan Neutrality Analysis}

Housing character weights satisfy the Shaw-Miller neutrality requirement
for the same structural reasons as all M-track signals: no electoral data,
voter registration data, or racial composition data enters any component of
the housing character vector.
The inputs are:
\begin{itemize}
  \item ACS B25024 (housing units by structural type) --- survey data,
    no racial or partisan content.
  \item ACS B25003 (tenure) --- survey data, no racial or partisan content.
  \item ACS B25035 (median year built) --- survey data, no racial or partisan content.
\end{itemize}
The correlation between housing character and race is incidental.
Single-family homeowner neighborhoods are documented to be disproportionately
white in many states~\citep{rothstein2017}, but the algorithm does not use
race; it uses the housing characteristics that happen to correlate with race
due to historical factors outside the algorithm's scope.
Under \textit{Shaw v.\ Reno}~\citeyearpar{shaw1993} and \textit{Miller v.\
Johnson}~\citeyearpar{miller1995}, the relevant question is whether race was
the \emph{predominant} factor.
An algorithm that uses only census-bureau-published housing survey data and
applies a deterministic cosine similarity formula cannot be characterized as
race-predominant redistricting.

\subsection{Expert Witness Deployment Template}

A redistricting practitioner deploying M.3 housing character weights should
structure expert testimony around the following elements:

\paragraph{1. Data provenance.}
``I used three tables from the U.S.\ Census Bureau's American Community Survey
5-Year Estimates: B25024 (units in structure), B25003 (tenure), and B25035
(median year structure built), all at census tract geography.
These are federal survey tables published at
\url{https://www.census.gov/programs-surveys/acs}.
No voter registration, election results, or racial composition data was used.''

\paragraph{2. Criterion alignment.}
``The [state]'s redistricting criteria require preservation of communities of
interest, which I interpret to include homeowner communities with shared
interests in neighborhood stability, historic preservation, and local school
quality.
I measured housing character as a four-component vector: the fraction of
single-family units, the fraction of large-multifamily units, the
owner-occupancy rate, and a vintage score reflecting the age of the housing stock.
District boundaries are drawn preferentially at the boundaries between tracts
with dissimilar housing character.''

\paragraph{3. Specific community preservation.}
``In the proposed plan, [specific historic neighborhood, e.g., `Cameron Park and
Five Points in Raleigh'] are placed within a single congressional district.
Both neighborhoods have median year built before 1955 and owner-occupancy rates
above 65\%.
Their housing character similarity score is [computed from the ACS housing
columns].
In the challenger's plan, these neighborhoods are split across two districts,
separating a community whose shared investment in historic preservation,
property values, and school quality has been documented over decades.''

\paragraph{4. When not to use M.3.}
\begin{itemize}
  \item States without a communities-of-interest criterion.
  \item States with uniform housing stock (no meaningful within-state housing
    character variation).
  \item Contexts where housing character weights would reduce minority
    opportunity districts below VRA thresholds.
\end{itemize}
